\documentclass[12pt,a4paper]{article}

\usepackage[utf8]{inputenc}
\usepackage[T1]{fontenc}
\usepackage[spanish]{babel}
\usepackage{geometry}
\usepackage{hyperref}
\usepackage{enumitem}

\geometry{margin=2.5cm}

\title{Manual de Usuario\Club55}
\author{}
\date{\today}

\begin{document}

\maketitle

\tableofcontents
\newpage

\section{Introducción}

Club55 es una plataforma diseñada para facilitar la participación de adultos mayores en actividades recreativas, culturales, educativas y de bienestar. Además, permite la interacción entre adultos mayores, cuidadores y proveedores de actividades.

\section{Acceso al sistema}

\subsection{Inicio de sesión}

Para ingresar al sistema:

\begin{enumerate}
\item Acceda a la página principal.
\item Seleccione la opción \textbf{Iniciar sesión}.
\item Ingrese una cuenta válida.
\item Presione \textbf{Entrar a mi cuenta}.
\end{enumerate}

\subsection{Cuentas de prueba}

\begin{itemize}
\item Adulto mayor:
\begin{itemize}
\item Correo: \texttt{[maria@club55.demo](mailto:maria@club55.demo)}
\end{itemize}

```
\item Cuidador:
\begin{itemize}
    \item Correo: \texttt{laura@club55.demo}
\end{itemize}

\item Proveedor:
\begin{itemize}
    \item Correo: \texttt{tertulia@club55.demo}
\end{itemize}
```

\end{itemize}

\section{Registro de usuario}

Para crear una cuenta nueva:

\begin{enumerate}
\item Acceda a la opción \textbf{Crear cuenta}.
\item Complete los datos solicitados:
\begin{itemize}
\item Nombre completo.
\item Correo electrónico.
\item Teléfono.
\item Ciudad.
\item Tipo de usuario.
\end{itemize}
\item Presione \textbf{Registrar}.
\end{enumerate}

\section{Exploración de actividades}

La sección \textbf{Actividades} permite visualizar todas las actividades disponibles.

Desde esta sección es posible:

\begin{itemize}
\item Consultar actividades disponibles.
\item Ver información detallada.
\item Revisar fecha, hora y ubicación.
\item Identificar la categoría de cada actividad.
\end{itemize}

\section{Reserva de actividades}

Para realizar una reserva:

\begin{enumerate}
\item Seleccione una actividad.
\item Revise la información disponible.
\item Presione el botón \textbf{Reservar}.
\item Confirme la operación.
\end{enumerate}

Una vez completado el proceso, la actividad aparecerá en la sección \textbf{Mis Reservas}.

\section{Gestión de reservas}

\subsection{Visualización de reservas}

La sección \textbf{Mis Reservas} permite consultar:

\begin{itemize}
\item Reservas activas.
\item Historial de reservas.
\item Reservas asociadas a un cuidador o adulto mayor vinculado.
\end{itemize}

\subsection{Filtros}

Las reservas pueden filtrarse por:

\begin{itemize}
\item Categoría.
\item Proveedor.
\item Fecha inicial.
\item Fecha final.
\end{itemize}

\subsection{Cancelación de reservas}

Para cancelar una reserva:

\begin{enumerate}
\item Ingrese a \textbf{Mis Reservas}.
\item Localice la actividad correspondiente.
\item Presione \textbf{Cancelar}.
\end{enumerate}

La reserva cambiará su estado y quedará registrada en el historial.

\section{Lista de espera}

Cuando una actividad alcanza su capacidad máxima, el usuario puede unirse a una lista de espera.

Si posteriormente se libera un cupo, el sistema podrá asignar la reserva automáticamente.

\section{Notificaciones}

La plataforma genera notificaciones para informar sobre:

\begin{itemize}
\item Confirmaciones de reserva.
\item Cancelaciones.
\item Cambios en actividades.
\item Avisos importantes del sistema.
\end{itemize}

\section{Mensajería}

La sección de mensajes permite la comunicación entre usuarios relacionados con las actividades.

Las funciones disponibles incluyen:

\begin{itemize}
\item Envío de mensajes.
\item Consulta de conversaciones.
\item Seguimiento de comunicaciones.
\end{itemize}

\section{Perfil de usuario}

Desde el perfil es posible:

\begin{itemize}
\item Modificar información personal.
\item Actualizar datos de contacto.
\item Revisar información de la cuenta.
\end{itemize}

\section{Funciones del proveedor}

Los usuarios con rol de proveedor pueden:

\begin{itemize}
\item Crear actividades.
\item Editar actividades existentes.
\item Cancelar actividades.
\item Consultar participantes inscritos.
\item Gestionar la oferta de actividades.
\end{itemize}

\section{Funciones del cuidador}

Los cuidadores pueden:

\begin{itemize}
\item Consultar actividades del adulto mayor vinculado.
\item Revisar reservas activas.
\item Consultar el historial de participación.
\end{itemize}

\section{Salud y bienestar}

La plataforma incluye contenido orientado al bienestar de los usuarios, como:

\begin{itemize}
\item Recomendaciones de salud.
\item Consejos de actividad física.
\item Hábitos de alimentación saludable.
\item Estrategias para el bienestar emocional.
\end{itemize}

\section{Juegos y entretenimiento}

La sección de juegos ofrece actividades recreativas orientadas al entretenimiento y estimulación cognitiva.

\section{Cierre de sesión}

Para finalizar la sesión:

\begin{enumerate}
\item Acceda al menú de usuario.
\item Seleccione la opción \textbf{Cerrar sesión}.
\end{enumerate}

\section{Conclusión}

Club55 permite a los adultos mayores participar en actividades de manera sencilla, mientras cuidadores y proveedores pueden gestionar y supervisar la participación mediante una interfaz accesible y centralizada.

\end{document}
